% !TEX program = xelatex
% ---------------------------------------------------------------------------
% 這份 .tex 是 resume pipeline 從 Notion 自動產生的，可以直接在 Overleaf 改。
%
% 1. 編譯器要用 XeLaTeX（不是 pdfLaTeX）—— 第一行的 magic comment 通常會自動選好；
%    沒有的話到 Menu → Compiler 改。xeCJK 與 fontspec 都需要它。
% 2. 字型是編譯時才決定的：找不到 Noto 就會自動退到 TeX Live 內建字型。
%    若仍因缺字型編不過，只要改下面那幾行 fallback 的字型名字就行。
% 3. 在這裡的手改只存在於 Overleaf。下一次 pipeline 產出會覆蓋這個檔案 ——
%    要永久生效請改 repo 裡的 templates/resume.tex.j2。
% ---------------------------------------------------------------------------
\documentclass[11pt,a4paper]{article}

\usepackage{fontspec}
\usepackage{xeCJK}

\IfFontExistsTF{Noto Serif}
  {\setmainfont{Noto Serif}}
  {\setmainfont{TeX Gyre Termes}}
\IfFontExistsTF{Noto Sans}
  {\setsansfont{Noto Sans}}
  {\setsansfont{TeX Gyre Heros}}
\IfFontExistsTF{Noto Serif CJK TC}
  {\setCJKmainfont{Noto Serif CJK TC}}
  {\IfFontExistsTF{uming.ttc}
    {\setCJKmainfont{uming.ttc}}
    {\setCJKmainfont{FandolSong-Regular.otf}}}
\IfFontExistsTF{Noto Sans CJK TC}
  {\setCJKsansfont{Noto Sans CJK TC}}
  {\IfFontExistsTF{ukai.ttc}
    {\setCJKsansfont{ukai.ttc}}
    {\setCJKsansfont{FandolHei-Regular.otf}}}

\IfFontExistsTF{Noto Sans CJK TC}
  {%
    \newfontfamily\symbolfont{Noto Sans CJK TC}
    \newcommand{\bulletmark}{{\symbolfont ▪}}
    \newcommand{\arrowright}{{\symbolfont →}}
    \newcommand{\arrowboth}{{\symbolfont ↔}}
  }
  {%
    \newcommand{\bulletmark}{\rule[0.05ex]{0.32em}{0.32em}}
    \newcommand{\arrowright}{\ensuremath{\rightarrow}}
    \newcommand{\arrowboth}{\ensuremath{\leftrightarrow}}
  }

\xeCJKDeclareCharClass{HalfLeft}{`—,`–}

\XeTeXlinebreaklocale "zh"
\XeTeXlinebreakskip = 0pt plus 1pt minus 0.1pt

\usepackage[
  a4paper,
  top=12mm,
  bottom=12mm,
  left=14mm,
  right=14mm
]{geometry}
\usepackage{titlesec}
\usepackage{enumitem}
\usepackage{xcolor}
\usepackage[hidelinks]{hyperref}
\usepackage{tabularx}
\usepackage{ragged2e}

\definecolor{ink}{HTML}{111111}
\definecolor{inkmuted}{HTML}{444444}
\color{ink}

\RaggedRight
\setlength{\parindent}{0pt}
\pagestyle{empty}

\titleformat{\section}
  {\normalfont\sffamily\bfseries\fontsize{10.5pt}{13pt}\selectfont\MakeUppercase}
  {}{0em}{}[{\vspace{-1.1mm}\rule{\linewidth}{0.5pt}}]
\titlespacing*{\section}{0pt}{4.2mm}{1.1mm}

\setlist[itemize]{
  leftmargin=3.4mm,
  itemsep=0.7mm,
  parsep=0pt,
  topsep=1mm,
  partopsep=0pt,
  label={\footnotesize\raisebox{0.15ex}{\bulletmark}}
}

\newcommand{\entryhead}[2]{%
  \begin{tabularx}{\linewidth}{@{}Xr@{}}
    \textbf{\fontsize{10.5pt}{13pt}\selectfont #1} &
    {\color{inkmuted}\fontsize{9.5pt}{12pt}\selectfont #2}
  \end{tabularx}%
}

\newcommand{\entrysub}[2]{%
  \begin{tabularx}{\linewidth}{@{}Xr@{}}
    {\color{inkmuted}\itshape \fontsize{9.5pt}{12pt}\selectfont #1} &
    {\color{inkmuted}\fontsize{9.5pt}{12pt}\selectfont #2}
  \end{tabularx}%
}

\fontsize{9.5pt}{13pt}\selectfont

\begin{document}

\begin{center}
  {\sffamily\bfseries\fontsize{17pt}{20pt}\selectfont
   WU, YU-HSUAN}\\[1.2mm]
  {\color{inkmuted}\fontsize{9pt}{11pt}\selectfont
   applegirlmabel@gmail.com}
\end{center}

\section{Education}
\entryhead{National Yang Ming Chiao Tung University (NYCU)}{2028 (expected)}
\entrysub{Master of Arts in Linguistics}{}
\begin{itemize}
  \item Relevant Coursework (expected): Computational Linguistics
\end{itemize}
\vspace{2.6mm}\entryhead{National Taiwan University (NTU)}{2026}
\entrysub{Bachelor of Arts in Foreign Languages and Literatures}{}
\begin{itemize}
  \item Member of NTU Computational Linguistics Specialized Program
  \item Relevant Coursework: NLP, Computational Linguistics, Phonology, Semantics
\end{itemize}

\section{Experience}
\entryhead{Cathay Financial Holdings — DDT AI}{Taipei, Taiwan}
\entrysub{AI Engineer \& System Design Intern}{Feb 2026 – Present}
\begin{itemize}
  \item Developed LLM agents for optimizing internal business use .Owning both the AI
application layer and the cloud-oriented deployment infrastructure.
\end{itemize}
\vspace{2.6mm}\entryhead{Taiwan AI Academy (AIA)}{New Taipei City, Taiwan}
\entrysub{AI Developer \& Technical Support Summer Intern}{Jun 2025 – Aug 2025}
\begin{itemize}
  \item A dual delivery-and-support role: built an internal-facing application while serving as front-line technical support for enterprise course participants.
\end{itemize}

\section{Projects}
\entryhead{Voice-Notion — Cloud-First Standup Assistant}{Cathay Financial Holdings, DDT AI}
\begin{itemize}
  \item Built a Slack bot that converts mp4 oral work report into confirmable Notion pages and tasks, plus a daily automated standup summary.
  \item Re-architected a single-instance Next.js service into a cloud-native ECS Fargate deployment with DynamoDB, SQS and EventBridge to enable horizontal scaling and zero-data-loss rolling updates.
  \item Evaluated Lambda + API Gateway, self-managed EC2 and App Runner against ECS Fargate on ops cost and private-subnet criteria.
\end{itemize}
\vspace{2.6mm}\entryhead{AI News Agent — Three-Scenario Routing Chatbot}{Cathay Financial Holdings, DDT AI}
\begin{itemize}
  \item Built an internal AI research report query assistant using a four-layer architecture (Interface / Application / Data / External Services).
  \item Designed a three-scenario classification system (LLM intent detection + Python state machine) for intelligent routing across explicit, fuzzy, and out-of-scope queries.
  \item Built a full ETL pipeline (Google Sheets \arrowright{} Bedrock Embedding \arrowright{} OpenSearch) with multi-dimensional filtering by topic, source, date, and subscription tier.
\end{itemize}
\vspace{2.6mm}\entryhead{BU Agent — Web-Based Two-Mode BRD Co-Drafting Assistant}{Cathay Financial Holdings, DDT AI}
\begin{itemize}
  \item Built an internal web app using LangGraph + PostgresSaver with a split layout (chat \arrowboth{} mode-aware workspace).
  \item BU Agent serves as the BU's single entry point (two modes) to draft v1.0 BRD, then hands off structured deliverables to the downstream BA Agent.
\end{itemize}

\section{Publications \& Talks}
\entryhead{Examining LLM's Ability to Judge Presuppositions: A Comparative Study of Zero-Shot and Self-Discover-NLI Prompting}{ROCLING 2025}
\begin{itemize}
  \item Developed and presented the \textquotedbl{}Self-Discover-NLI\textquotedbl{} prompt engineering framework, adapting the three-step Self-Discover pipeline (SELECT \arrowright{} ADAPT \arrowright{} APPLY) specifically for Natural Language Inference (NLI) tasks.
  \item Conducted systematic experiments on the IMPPRES dataset (95 tasks per trigger type), evaluating GPT-4o's handling of clefts, conditionals, modals, and other linguistic triggers.
  \item Achieved \textasciitilde{}10\% overall accuracy improvement; performed in-depth error analysis identifying systematic model weaknesses.
\end{itemize}

\section{Skills}
\begin{itemize}[label={},leftmargin=0pt,itemsep=0.7mm]
  \item \textbf{Languages:}\ Mandarin (Native), English (Fluent)
  \item \textbf{Programming:}\ Python
  \item \textbf{Frameworks:}\ LangGraph, LangChain
  \item \textbf{Tools:}\ AWS CLI, Docker
  \item \textbf{Certificates:}\ TOEIC (Listening and Reading 985, Speaking 180, Writing 170), AWS SAA (Solutions Architect Associate)
\end{itemize}

\end{document}